\chapter{Security: JWT, OAuth2, and the Trusted Edge}
\label{ch:security}

Security is where this project is most instructive, because it uses three
mechanisms that industry uses constantly --- stateless JWTs, OAuth2 login,
TOTP two-factor --- and wires them through a gateway in the standard
``trusted edge'' pattern.

\section{Spring Security in one page}

Spring Security is a chain of servlet filters in front of your
controllers. Each request passes through the chain; filters authenticate
(who are you?) and the authorization rules decide (may you?). The
project's auth-service chain, from its \code{SecurityConfig}:

\begin{lstlisting}[language=Java]
http
  .csrf(AbstractHttpConfigurer::disable)                    (*@\co{1}@*)
  .sessionManagement(s ->
      s.sessionCreationPolicy(SessionCreationPolicy.STATELESS)) (*@\co{2}@*)
  .authorizeHttpRequests(auth -> auth
      .requestMatchers("/register", "/login", "/activate",
          "/refresh", "/verify-2fa", "/change-password",
          "/oauth2/**", "/login/oauth2/**",
          "/v3/api-docs/**", "/swagger-ui/**").permitAll()  (*@\co{3}@*)
      .anyRequest().authenticated())
  .oauth2Login(o -> o.successHandler(oAuth2LoginSuccessHandler))
  .addFilterBefore(gatewayHeaderAuthFilter,
      UsernamePasswordAuthenticationFilter.class);          (*@\co{4}@*)
\end{lstlisting}

\begin{description}
  \item[\co{1}] CSRF protection defends browser cookie sessions; a
  stateless JWT API has no cookie session to ride, so it is disabled.
  Disable it for that reason --- not because a tutorial did.
  \item[\co{2}] \code{STATELESS}: no HTTP session is created. Every
  request must carry its own proof of identity. This is the property that
  lets any replica answer any request --- the scaling prerequisite.
  \item[\co{3}] The public paths. Note they are the \emph{stripped} paths
  (\code{/login}, not \code{/api/auth/login}) --- Chapter~4 explains why.
  Note also what is \emph{absent}: \code{/actuator/**}. That absence
  decides this service's Kubernetes probe strategy in Chapter~17.
  \item[\co{4}] The custom filter that trusts identity headers from the
  gateway --- the subject of Section~\ref{sec:trusted-edge}.
\end{description}

\section{JWT: stateless identity}

A JSON Web Token is three base64url parts, \code{header.payload.signature}.
The payload carries \emph{claims}: subject (user id), role, expiry. The
signature is HMAC-SHA256 over header and payload with a shared secret ---
the 256-bit-minimum string both auth-service and the gateway hold.

The properties that matter:

\begin{itemize}
  \item \textbf{Self-contained.} Validation needs no database call: check
  signature, check expiry, read claims. This is what lets the
  \emph{gateway} authenticate without owning user data.
  \item \textbf{Readable.} Base64 is encoding, not encryption --- anyone
  can decode the payload. Never put secrets in claims.
  \item \textbf{Irrevocable until expiry.} A signed token cannot be
  un-signed. Hence the two-token pattern: a short-lived \emph{access
  token} (here 15 minutes) paired with a long-lived \emph{refresh token}
  (7 days) that \emph{is} checked against the database at
  \code{/refresh} --- revocation happens there.
\end{itemize}

\section{The trusted-edge pattern}
\label{sec:trusted-edge}

The gateway's \code{JwtAuthenticationFilter} validates the token, then
forwards the request with identity as plain headers:

\begin{lstlisting}[language=Java]
ServerHttpRequest mutatedRequest = exchange.getRequest().mutate()
    .header(JwtConstants.HEADER_USER_ID, userId)
    .header(JwtConstants.HEADER_USER_ROLE, role)
    .build();
\end{lstlisting}

Downstream, each service's \code{GatewayHeaderAuthFilter} reads those
headers and builds the Spring Security context from them --- no
re-validation. Services stay simple; cryptography happens once.

The pattern's obligation: those headers are only trustworthy if
\emph{every} path to a service goes through the gateway. A client who can
reach course-service directly could forge \code{X-User-Id}. This is the
deep reason the direct \code{/courses} Ingress was deleted (Chapter~4),
and why in-cluster traffic is the remaining soft spot --- any pod can
still call \code{course-service:8082} with invented headers.
NetworkPolicies or mTLS (service mesh) are the industrial hardening;
Chapter~26 discusses when they earn their complexity.

\section{OAuth2 login with Google}

OAuth2's \emph{authorization code} flow, as it runs here:

\begin{enumerate}
  \item Browser hits \code{/oauth2/authorization/google} (via the
  gateway); Spring redirects to Google with the app's client-id and a
  \code{redirect\_uri}.
  \item The user consents at Google; Google redirects the browser to the
  \code{redirect\_uri} --- \code{/login/oauth2/code/google} --- carrying a
  one-time \emph{code}.
  \item Auth-service exchanges code + client-secret for tokens
  server-to-server, reads the user's e-mail/profile, then the project's
  \code{OAuth2LoginSuccessHandler} takes over: find or create the local
  account and issue \emph{our own} JWTs. Google authenticates; the
  application still authorizes with its own tokens.
\end{enumerate}

\begin{pitfall}[redirect\_uri\_mismatch]
Symptom: Google shows ``Error 400: redirect\_uri\_mismatch'' instead of a
consent screen. Cause: the URI the app sent is not byte-identical to one
registered in the Google Cloud console. It changes per environment
(\code{localhost:8080} in dev, \code{ts.local} in the cluster --- which is
why the Deployment overrides
\code{...GOOGLE\_REDIRECT\_URI}), and every variant must be registered.
The URI must reach auth-service \emph{unstripped}, which is why the
OAuth2 gateway routes carry no StripPrefix.
\end{pitfall}

\section{TOTP two-factor}

Time-based One-Time Passwords (RFC~6238): at enrollment the server
generates a secret, shows it as a QR code, and the authenticator app
derives a 6-digit code from \code{HMAC(secret, time/30s)}. Login becomes
two steps --- password check issues a short-lived pending state, then
\code{/verify-2fa} checks the code and issues the real tokens. Server and
phone never communicate; they only share the secret and a clock.

\section{Outside the project}

The main industrial alternative moves token issuing out of your code
entirely: an identity provider (Keycloak self-hosted, or Auth0/Cognito
managed) issues tokens signed with an \emph{asymmetric} keypair
(RS256). Services validate with the \emph{public} key fetched from a
JWKS URL --- no shared secret to distribute, key rotation built in:

\begin{lstlisting}[language=yaml]
spring:
  security:
    oauth2:
      resourceserver:
        jwt:
          jwk-set-uri: https://idp.company.com/realms/app/protocol/openid-connect/certs
\end{lstlisting}

Understanding the hand-rolled HS256 version --- what this project builds
--- is precisely what makes that configuration meaningful: same claims,
same validation, different key ceremony.

\begin{quiz}
\begin{enumerate}
  \item Why can CSRF protection be disabled for a stateless JWT API but
  not for a cookie-session webapp?
  \item An attacker steals an access token. How long is it useful, what
  cannot the server do about it meanwhile, and which endpoint finally
  stops the session?
  \item In the trusted-edge pattern, state the exact invariant that must
  hold for \code{X-User-Id} to be trustworthy, and name two deployment
  mistakes that would break it.
  \item In the Google flow, which secret never leaves the server, which
  value transits the browser, and why is the browser-visible one safe?
  \item HS256 vs RS256: who can validate, who can sign, and why do
  identity providers choose RS256?
\end{enumerate}
\end{quiz}
